% Gauss, Werke X/1 (1917), native128 lower–133, printed121–126.
\NumberedItem{2}
\begin{center}\large\textsc{Gauss} an \textsc{Gerling}.\end{center}
\noindent\hspace*{\parindent}Sie haben wol gezürnt, liebster \textsc{Gerling}, dass ich Ihr Verlangen,
Ihre Abhandlung Ihnen noch im vorigen Jahre zurückzusenden, nicht erfüllt habe.
Theils wünschte ich Ihnen über die Polygone die nöthige Auskunft zu geben; allein immer
konnte ich dazu keine Musse gewinnen, indem besonders die practischen Arbeiten fast alle
Zeit absorbiren. Theils hoffte ich, Ihnen ein
\par\vfill\noindent X1.\hfill16\par
\clearpage
\SourceLeaf{129}\ArithmeticHead{122}{L}{ALGEBRA.\quad NACHLASS.}
\noindent Exemplar meiner Vorlesung über die Attraction elliptischer Ringe [*)] beifügen
zu können, welches auch manches andre Ihnen wol Interessante enthält z. B. die ersten Linien
meiner arithmetisch-geometrischen Mittel; allein obgleich ich diese Vorlesung bereits vor
einigen Wochen zur Correctur gehabt, habe ich doch noch immer keine Abdrücke erhalten.
Länger darf ich aber mit meiner Antwort nicht zögern.

Zuvörderst danke ich verbindlichst für die Mittheilung Ihres Aufsatzes, den ich mit vielem
Vergnügen gelesen habe. Ihre Resultate sind merkwürdig und ich wünsche sehr, dass Sie diese
Versuche weiter fortsetzen mögen.

Ich glaube, dass Sie mit Unrecht die Theorie der Polygone in meinen D. A. scheuen,
und dass es nur auf den Versuch ankommt, wo Sie gewiss alles sehr leicht finden werden.
Wenn ich jezt ein Paar Stunden dazu anwende, über diesen Gegenstand zu schreiben,
so geschieht es in der Hoffnung, dass eben das Lückenhafte, welches in einer abgerissenen
Darstellung der 17-Theilung, die ich hier gebe, nothwendig bleiben muss, Sie anreizen wird,
das Allgemeine an der Quelle zu schöpfen. Für mich wenigstens sind und bleiben die
Unters[uchungen] der höhern Arithmetik, bei weitem das Allerschönste der Mathematik,
und der Genuss, den ich, auch an der schönsten astronomischen Untersuchung finde,
ist gar Nichts, verglichen mit dem, welchen die höhere Arithmetik gewährt.

Bloss synthetisch beweisen lässt sich die Constructibilität des 17-Ecks sehr leicht und kurz.

Es sei
\[360^0=17\varphi.\]
\noindent Ich setze
\[\begin{array}{l|l}
\cos\varphi+\cos4\varphi=a&a+b=e\\
\cos2\varphi+\cos8\varphi=b&c+d=f\\
\cos3\varphi+\cos5\varphi=c&\text{Also nach einem bekannten Satze}\\
\cos6\varphi+\cos7\varphi=d&1)\ e+f=-\tfrac12.
\end{array}\]
\noindent Durch leichte Entwickelung, mit der Überlegung, dass
\[\cos n\varphi=\cos(17-n)\varphi,\]
\par\smallskip\noindent\rule{.25\linewidth}{.4pt}\par\smallskip
{\small\noindent\hspace*{\parindent}[*) \textit{Determinatio attractionis etc.,} 1818,
Werke III, S. 331.]\par}
\clearpage
\SourceLeaf{130}\ArithmeticHead{123}{R}{ÜBER DAS REGELMÄSSIGE SIEBZEHNECK.}
\noindent findet man
\[\begin{array}{l|l}
2ab=e+f=-\tfrac12&2bc=a+2c+d\\
2ac=2a+b+d&2bd=a+2b+c\\
2ad=b+c+2d&2cd=e+f=-\tfrac12.
\end{array}\]
\noindent Hieraus also
\[2ac+2ad+2bc+2bd=4a+4b+4c+4d\]
\noindent d. i.
\[2ef=-2\]
\noindent oder
\[\makebox[0pt][r]{\hspace*{-10em}2)\hspace{8em}}ef=-1.\]
\noindent\hspace*{\parindent}Man sieht also aus 1) und 2), dass $e$ und $f$ Wurzeln der Gleichung
\[xx+\tfrac12x-1=0\]
\noindent sind also die eine gleich
\[=-\tfrac14+\surd\tfrac{17}{16},\ \text{die andre }=-\tfrac14-\surd\tfrac{17}{16}.\]
\noindent Dass die erste $=e$, die andre $=f$ ist, gibt eine oberflächliche Kenntniss der
numerischen Werthe; will man diese nicht benutzen, so ist die Untersuchung altioris indaginis
(besonders bei den Polygonen im \textls[120]{Allgemeinen}) und kann hier nicht entwickelt werden.
S. m. Abhandlung \textit{Summatio quarundam serierum singularium} 1808 [*)].

Ferner sind $a$ und $b$ Wurzeln der Gleichung
\[xx-ex-\tfrac14=0\]
\noindent also die Werthe
\[\begin{aligned}
&\tfrac12e\pm\surd(\tfrac14+\tfrac14ee)\\
&=-\tfrac18+\tfrac18\surd17\pm\tfrac18\surd(34-2\surd17).
\end{aligned}\]
\noindent Dass das obere Zeichen für $a$, das untere für $b$ gelten muss, ist in
\textls[120]{diesem} Falle leicht einzusehen, da
\[a-b=(\cos\varphi-\cos2\varphi)+(\cos4\varphi-\cos8\varphi)\]
\noindent offenbar positiv seyn muss. Auf ähnliche Weise wird
\par\smallskip\noindent\rule{.25\linewidth}{.4pt}\par\smallskip
{\small\noindent\hspace*{\parindent}[*) Werke II, S. 9.]\par}
\par\vfill\noindent\hfill16\textsuperscript{*}\par
\clearpage
\SourceLeaf{131}\ArithmeticHead{124}{L}{ALGEBRA.\quad NACHLASS.}
\[\begin{aligned}
c&=-\tfrac18-\tfrac18\surd17+\tfrac18\surd(34+2\surd17)\\
d&=-\tfrac18-\tfrac18\surd17-\tfrac18\surd(34+2\surd17).
\end{aligned}\]
\noindent Endlich sind nun offenbar $\cos\varphi$ und $\cos4\varphi$ die Wurzeln der quadratischen
Gleichung (weil das Product $\cos\varphi\ \cos4\varphi=\tfrac12c$)
\[xx-ax+\tfrac12c\ [=0],\]
\noindent also
\[\begin{aligned}
\cos\varphi&=+\tfrac12a+\surd(\tfrac14aa-\tfrac12c)\\
\cos4\varphi&=+\tfrac12a-\surd(\tfrac14aa-\tfrac12c).
\end{aligned}\]
\noindent Da nun
\[2aa=2+b+2c\]
\noindent wird, so ist
\[\begin{aligned}
\cos\varphi&=\tfrac12a+\surd(\tfrac14+\tfrac18b-\tfrac14c)\\
&=-\tfrac1{16}+\tfrac1{16}\surd17+\tfrac1{16}\surd(34-2\surd17)
\overset{*}{+}\tfrac18\surd\{17+3\surd17\\
&\qquad\qquad{}-\surd(34-2\surd17)-2\surd(34+2\surd17)\}.
\end{aligned}\]
\noindent Diess ist dieselbe Formel, die in meinen D. A. p. 662 steht [*)], nur ist dort
durch einen Druckfehler statt des $+$, welches hier mit $*$ bezeichnet ist, ein $-$ gesetzt,
oder, was dasselbe ist, die dortige Formel stellt nicht $\cos\varphi$ sondern $\cos4\varphi$
d. i. $\sin\dfrac{90^0}{17}$ vor, also die halbe Seite des 34-Ecks. — Die übrigen Cosinus
bekommen ganz ähnliche Ausdrücke.

In den D. A. habe ich vorgezogen, statt des $\cos n\varphi$, die Grösse
\[\cos n\varphi+\surd-1\sin n\varphi\]
\noindent d. i. die Wurzel der Gleichung $x^{17}-1=0$ zu brauchen, wodurch alles viel
zierlicher wird.

Die obige Darstellung ist an sich zureichend. Im Allgemeinen, wenn es die Theilung des
Kreises in $p$ Theile gilt, wo $p$ eine Primzahl, kommt es darauf an, die
$\dfrac{p-1}{2}$ Cosinus der Bögen
\[\frac{360}p,\ 2\mathpunct{.}\frac{360}p,\ 3\mathpunct{.}\frac{360}p,\ldots,
\frac12(p-1)\mathpunct{.}\frac{360}p\]
\par\smallskip\noindent\rule{.25\linewidth}{.4pt}\par\smallskip
{\small\noindent\hspace*{\parindent}[*) Werke I, S. 462, wo der Druckfehler verbessert ist.]\par}
\clearpage
\SourceLeaf{132}\ArithmeticHead{125}{R}{ÜBER DAS REGELMÄSSIGE SIEBZEHNECK.}
\noindent auf eine solche Weise in Gruppen zu theilen, dass ähnliche Erfolge hervorgehn wie hier,
und diess beruht lediglich [auf] arithmetischen Principien, die der 3. Abschnitt meiner
\textit{Disq} enthält. Alles hängt dabei von den Factoren der Zahl $\tfrac12(p-1)$ ab;
ist diese Zahl eine Potenz von 2, z. B.
\[p=3,\ 5,\ 17,\ 257,\ 65\,537,\]
\noindent so kommen bloss quadratische Gleichungen vor; hingegen z. B. für $p=31$, wo
$\tfrac12(p-1)=3\mathpunct{.}5$, ist eine cubische und eine Gleichung vom 5. Grade unausweichlich.

Das Geschichtliche jener Entdeckung ist bisher nirgends von mir öffentl. erwähnt;
ich kann es aber sehr genau angeben. Der Tag war der 29. März 1796, und der Zufall hatte
gar keinen Antheil daran. Schon früher war alles was auf die Zertheilung der Wurzeln der Gleichung
\[\frac{x^p-1}{x-1}=0\]
\noindent in zwei Gruppen sich bezieht, von mir gefunden, wovon der schöne Lehrsatz D. A. p. 637 [*)]
unten abhängt u. zwar im Winter 1796 (meinem ersten Semester in Göttingen), ohne dass ich den Tag
aufgezeichnet hätte. Durch angestrengtes Nachdenken über den Zusammenhang aller Wurzeln unter
einander nach arithmetischen Gründen, glückte es mir bei einem Ferienaufenthalt in Braunschweig,
am Morgen des gedachten Tages (ehe ich aus dem Bette aufgestanden war) diesen Zusammenhang auf
das klarste anzuschauen, so dass ich die specielle Anwendung auf das 17-Eck und die numerische
Bestätigung auf der Stelle machen konnte. Freilich sind später noch andere Untersuchungen des
7. Abschnitts der D. A. hinzu gekommen. Ich kündigte diese Entdeckung in der Jenaischen
Literaturzeitung an, wo mein Inserat ungefähr im May oder Junius 1796 [**)] abgedruckt seyn wird.
Der Druck meiner \textit{Disq. Arith.} fing im April 1798 an, ging langsam fort, wurde mehrere male
ganz unterbrochen (weil der Drucker von Braunschweig wegzog, daher vom Bogen R an, andere Schrift
gebraucht ist) und wurde im Sommer 1801 vollendet.

Im Jahr 1798 oder 1799 kam ein gewisser \textsc{Huguené} oder \textsc{Huguenot} (ein
\par\smallskip\noindent\rule{.25\linewidth}{.4pt}\par\smallskip
{\small\noindent\hspace*{\parindent}[*) Werke I, S. 443.]\par
\noindent\hspace*{\parindent}[**) Siehe S. 3 dieses Bandes; die Anzeige ist in der That im Juni 1796 erschienen.]\par}
\clearpage
\SourceLeaf{133}\ArithmeticHead{126}{L}{ALGEBRA.\quad NACHLASS.}
\noindent preussischer Officier) durch Braunschweig, dem v. \textsc{Zimmermann} von meinen
Untersuchungen erzählte. Auf seine Bitte theilte ich ihm einen kleinen Aufsatz, der die Theorie
des 17-Ecks vollständig enthielt (ungefähr wie oben, aber viel ausführlicher) mit, den er behielt.
Dieser Mensch soll sich nachher nicht entblödet haben, in einem Werke, das er hat drucken lassen
[*)], welches ich aber selbst nicht gesehen habe, diese Sache auf eine solche Art vorzutragen,
dass der Leser glauben könnte, es sei des \textsc{Huguenot} eigne Arbeit! Unglücklicherweise aber,
soll dieses Buch sogar mehrere Jahre nach der Erscheinung meiner D. A. erst gedruckt seyn,
wodurch seine Unverschämtheit umso lächerlicher wird [**)]. — Übrigens darf ich Ihrer Discretion
und Ihrer Beurtheilung vertrauen, falls Sie von dem Geschichtlichen irgend etwas zu erwähnen
passend finden. .\,.\,.\,.\,.\,.\,.\,.\,.\,.\,.\,.\,.\,.\,.\,.\,.\,.\,.\,.\,.\,.\,.\,.

Leben Sie wohl, liebster \textsc{Gerling} und erfreuen bald mit einigen Zeilen
\begin{flushright}Ihren ergebensten\end{flushright}
\noindent\hspace*{\parindent}Göttingen d. 6. Januar 1819.\hfill C. F. G.\par
\medskip
Soeben, indem ich das Paket absenden will, werden mir einige Abdrücke meiner Vorlesung
überbracht; ich ergreife also diese Gelegenheit, noch einen für Sie beizulegen.
\par\smallskip\noindent\rule{.25\linewidth}{.4pt}\par\smallskip
{\small\noindent\hspace*{\parindent}[*) \textit{Mathematische Beyträge zur weiteren Ausbildung
angehender Geometer,} von dem K. Preuss. Hauptmann im Feld-Artillerie Corps v. \textsc{Huguenin},
Königsberg, bey Goebbels \& Unzer, 1803.]\par
\noindent\hspace*{\parindent}[**) \textsc{Gauss} war über das Buch v. \textsc{Huguenins} ungenau
unterrichtet. Auf S. 283 dieses Buches findet sich nämlich die folgende Fußnote:
»Dass das 17-Eck durch quadratische Gleichungen in einen Kreis beschrieben werden könne,
entdeckte zuerst ein junger Geometer aus Braunschweig, dessen Name wie ich mich erinnere
\textsc{Gaus} ist: er machte solches einigen Gelehrten bekannt und da dieses Verwundrung und
vielleicht einigen Unglauben erregte, so gab er etwas über sein Verfahren an. Zu dieser Zeit
(im J. 1796) befand ich mich in Braunschweig, wo ich solches nebst allem was man von dieser
unerwarteten Auflösung wusste, erfuhr und hatte das Vergnügen, diese Spur verfolgend,
seine \textls[120]{eigene Auflösung} in kurzer Zeit heraus zu bringen; die hier gegebene Auflösung
aber ist meine eigene und gänzlich von jener verschieden; daher es mir frey stehen wird,
diese als mein Eigenthum hier bekannt zu machen, indem ich die Herausgabe der ersten Auflösung,
welche gänzlich trigonometrisch war, ihrem Erfinder selbst überlasse«.]\par}
\ClosingRule
\clearpage
